\documentclass[11pt]{article}
\usepackage[a4paper,margin=1in]{geometry}
\usepackage{amsmath,amssymb,amsthm}
\usepackage{enumitem}
\usepackage{mathrsfs}
\usepackage{hyperref}

\DeclareMathOperator{\Spec}{Spec}
\DeclareMathOperator{\Max}{Max}
\DeclareMathOperator{\Rad}{Rad}
\DeclareMathOperator{\Nil}{Nil}
\DeclareMathOperator{\Frac}{Frac}
\newcommand{\fp}{\mathfrak p}
\newcommand{\fm}{\mathfrak m}
\newcommand{\fa}{\mathfrak a}
\newcommand{\fb}{\mathfrak b}
\newcommand{\fq}{\mathfrak q}
\newcommand{\fr}{\mathfrak r}
\newcommand{\fN}{\mathfrak N}

\title{Errata for \emph{Chapter 1 --- Rings and Ideals}}
\author{}
\date{\today}

\begin{document}
\maketitle

\noindent
Exercise numbers below refer to the order of the \texttt{exercise} environments
in the original file, which coincides with Atiyah--MacDonald, Chapter~1,
Exercises~1--28.  For each item, \textbf{Issue} describes the problem and
\textbf{Correction} gives a fix.  Exercises which are only sketches or are
unfinished are collected in the last two sections.

\section*{Summary of serious mathematical errors}

\begin{itemize}[leftmargin=1.5em]
  \item \textbf{Ex.~2(iii)(a)}: the displayed system of equations is wrong, and
        the conclusion \(f b_m^n=0\) does not produce a nonzero annihilator.
  \item \textbf{Ex.~2(iv)}: the proof of ``\(f,g\) primitive
        \(\Rightarrow fg\) primitive'' is not valid.
  \item \textbf{Ex.~4}: \(A[x]\) need not have Krull dimension \(0\); the
        asserted statement ``every prime ideal of \(A[x]\) is maximal'' is
        false.
  \item \textbf{Ex.~5(iii)(b)}: the claim that an arbitrary ideal equals the
        full preimage of its constant-term ideal is false, so the
        classification of maximal ideals is not proved.
  \item \textbf{Ex.~6}: the logical reduction is incorrect and the proof is
        unfinished.
  \item \textbf{Ex.~8}: the nilradical is not a prime ideal in general, so it
        cannot be used as a minimal prime.
  \item \textbf{Ex.~13}: the minimal-degree argument is not written correctly.
  \item \textbf{Ex.~14(ii)}: the proof that a maximal element of \(\Sigma\) is
        prime is invalid.
  \item \textbf{Ex.~16}: the zero ideal is missing from
        \(\Spec\mathbb Z\), \(\Spec\mathbb C[x]\), and
        \(\Spec\mathbb R[x]\).
  \item \textbf{Ex.~17(vi),(vii)}: the radical is confused with the ideal
        itself, and the quasi-compactness argument in (vii) is wrong.
  \item \textbf{Ex.~18(iv)}: \(V(\fp_x)\) is closed but is not a neighbourhood
        of \(x\) in general; the \(T_0\) proof is not valid as written.
  \item \textbf{Ex.~19}: the category-theoretic argument does not prove
        irreducibility of \(\Spec(A)\).
  \item \textbf{Ex.~21(v)}: the chain of equivalences is not justified; the
        correct criterion is \(\ker\phi\subseteq\Nil(A)\).
  \item \textbf{Ex.~23(iii),(iv)}: the proof that clopen sets are basic is
        wrong, and \(X\setminus\{q\}\) is not closed.
  \item \textbf{Ex.~24}: the associativity computation for addition and the
        absorption computation are incorrect.
  \item \textbf{Ex.~26}: the implication ``\(f\in\fm_y\Rightarrow f(y)\neq0\)''
        has the wrong sign.
  \item \textbf{Ex.~27}: the argument uses the weak Nullstellensatz, which is
        exactly what is not yet available at this point.
  \item \textbf{Ex.~28}: the crucial verification that the constructed map
        lands in \(Y\) is missing, and injectivity is not proved.
\end{itemize}

\section*{Exercise 1}

\textbf{Issue.}  The identity
\[
(1+x)(1-x+x^2-\cdots+(-1)^{r-1}x^{r-1})=1-x^r
\]
is not correct in general.  The correct identity is
\[
(1+x)(1-x+x^2-\cdots+(-1)^{r-1}x^{r-1})=1-(-x)^r.
\]
Since \(x^r=0\), both \(1-x^r\) and \(1-(-x)^r\) equal \(1\), so the final
conclusion is correct.  The exercise also asks to deduce that the sum of a unit
and a nilpotent element is a unit; this deduction is not written.

\textbf{Correction.}  Write
\[
(1+x)(1-x+\cdots+(-1)^{r-1}x^{r-1})=1-(-x)^r=1.
\]
For the deduction, if \(u\) is a unit and \(y^n=0\), then
\[
u+y=u(1+u^{-1}y),
\]
and \((u^{-1}y)^n=0\).  By the first part, \(1+u^{-1}y\) is a unit, hence
\(u+y\) is a unit.

\section*{Exercise 2}

\subsection*{(i)(a)}
Correct.  The displayed system
\[
a_nb_m=0,\quad a_n^2b_{m-1}=0,\quad\ldots,\quad a_n^{m+1}b_0=0
\]
is obtained by multiplying the coefficients of \(x^{n+m},x^{n+m-1},\ldots\)
by suitable powers of \(a_n\).  Since \(a_0b_0=1\), we get
\[
a_n^{m+1}=a_n^{m+1}b_0a_0=0.
\]
Then \(f-a_nx^n\) is a unit by Exercise~1, and iteration gives the
nilpotence of \(a_{n-1},\ldots,a_1\).

\subsection*{(i)(b)}
Correct, except for the typo \(a_ns^n\) in line~87 of the original; it should
be \(a_nx^n\).

\subsection*{(ii)(a)}
The idea is correct, but one step needs justification: from the nilpotence of
\[
xh=a_1x+a_2x^2+\cdots+a_nx^n
\]
one cannot immediately conclude that \(h=a_1+a_2x+\cdots+a_nx^{n-1}\) is
nilpotent without using that \(x\) is a non-zero-divisor in \(A[x]\).  Indeed,
if \((xh)^k=0\), then \(x^kh^k=0\), hence \(h^k=0\).

\subsection*{(ii)(b)}
The word ``obviously'' hides a short argument.  If each \(a_i\) is nilpotent,
then each \(a_ix^i\) is nilpotent, and a finite sum of nilpotent elements is
nilpotent.  Hence \(f\) is nilpotent.

\subsection*{(iii)(a): incorrect system of equations}

\textbf{Issue.}  The displayed system
\[
a_nb_m=0,\quad a_{n-1}b_m=0,\quad a_{n-2}b_m^2=0,\quad\ldots,\quad
a_1b_m^{n-1}=0,\quad a_0b_m^n=0
\]
is not correct.  For example, comparing the coefficient of \(x^{n+m-1}\) gives
\[
a_{n-1}b_m+a_nb_{m-1}=0.
\]
Multiplying by \(b_m\) and using \(a_nb_m=0\) gives
\[
a_{n-1}b_m^2=0,
\]
not \(a_{n-1}b_m=0\).  In general, without using minimality of \(m\), one only
gets
\[
a_{n-r}b_m^{r+1}=0\qquad(0\le r\le n),
\]
hence \(fb_m^{n+1}=0\).  But \(b_m\neq0\) does not imply \(b_m^{n+1}\neq0\), so
this does not give the required nonzero annihilator.

A concrete counterexample to the displayed system is the following.  Let
\[
A=k[x,y]/(x^2,y^2),\qquad f=x+yX,\qquad g=-x+yX.
\]
Then
\[
fg=-x^2+xyX-xyX+y^2X^2=0,
\]
but with \(a_0=x,a_1=y,b_0=-x,b_1=y\) we have
\[
a_0b_1=xy\neq0.
\]
Thus the equation \(a_0b_m^n=0\) is false without further hypotheses.

\textbf{Correction.}  Use the least-degree annihilator, as in the hint.  Let
\[
g=b_0+b_1x+\cdots+b_mx^m\neq0
\]
have minimal degree \(m\) among all nonzero polynomials with \(fg=0\).  Then
\(b_m\neq0\).  If \(m=0\), then \(g=b_0\) is a nonzero constant and
\(fb_0=0\), so we are done.  Assume \(m\ge1\).  Since the coefficient of
\(x^{n+m}\) is \(a_nb_m\), we have
\(a_nb_m=0\).  Hence \(a_ng\) has degree \(<m\) and still annihilates \(f\):
\[
f(a_ng)=a_n(fg)=0.
\]
By minimality of \(m\), \(a_ng=0\).  Comparing the coefficient of
\(x^{n+m-1}\) now gives \(a_{n-1}b_m=0\), so \(a_{n-1}g\) has degree \(<m\)
and annihilates \(f\); by minimality, \(a_{n-1}g=0\).  Continuing, we obtain
\[
a_{n-r}g=0\qquad(0\le r\le n).
\]
In particular \(a_ib_m=0\) for every \(i\), hence
\[
fb_m=0.
\]
Since \(b_m\neq0\), this is the desired nonzero element annihilating \(f\).

\subsection*{(iii)(b)}
Correct: if \(af=0\) with \(a\neq0\), then \(f\) is a zero-divisor by
definition.

\subsection*{(iv)(\(\Leftarrow\)): invalid proof of primitivity}

\textbf{Issue.}  In the proof that \(f,g\) primitive implies \(fg\) primitive,
the displayed formula beginning with ``\(fg=\)'' is not actually \(fg\).
It is the weighted sum
\[
\sum_{i,j}\lambda_i\mu_jb_ia_jx^{i+j}.
\]
Expressing \(1\) as a linear combination of the individual products
\(a_ib_j\) does not show that \(1\) lies in the ideal generated by the
\emph{coefficients} of \(fg\), which are the anti-diagonal sums
\[
c_k=\sum_{i+j=k}a_jb_i.
\]
The coefficients \(\lambda_i\mu_j\) are not constant on an anti-diagonal, so
the terms cannot be regrouped as an \(A\)-linear combination of the \(c_k\).

\textbf{Correction.}  Reduce modulo a maximal ideal.  Let \(\fm\) be a maximal
ideal of \(A\).  In the polynomial ring \((A/\fm)[x]\), which is a domain, the
image of \(fg\) is \(\bar f\bar g\).

If \(f,g\) are primitive and \(fg\) were not primitive, then there is a
maximal ideal \(\fm\) containing all coefficients of \(fg\), so
\(\bar f\bar g=0\).  Since \((A/\fm)[x]\) is a domain, \(\bar f=0\) or
\(\bar g=0\), contradicting the primitivity of \(f\) or \(g\).

Conversely, if \(fg\) is primitive and \(f\) were not primitive, then some
maximal ideal \(\fm\) contains all coefficients of \(f\); then
\(\bar f=0\), so \(\bar f\bar g=0\), contradicting the primitivity of \(fg\).
The same argument applies to \(g\).  Hence \(f\) and \(g\) are primitive.

The direction \(fg\) primitive \(\Rightarrow g\) primitive is also proved by
the rearrangement in the original text, provided the last term is corrected
from \(a_n+\xi_{m+n}\) to \(a_n\xi_{m+n}\).  The problematic direction is
the one discussed above.

\subsection*{Remark after Exercise 2}
The remark says ``\(n+m-1\) anti-diagonals''.  The coefficients of \(fg\) run
from \(x^0\) to \(x^{n+m}\), so there are \(n+m+1\) anti-diagonals, not
\(n+m-1\).

\section*{Exercise 3}

No mathematical error, but the induction is only indicated.  For example, for
(i) apply Exercise~2(i) with the coefficient ring
\(R=A[x_1,\ldots,x_{n-1}]\): a polynomial in \(x_n\) is a unit in \(R[x_n]\)
if and only if its constant term is a unit in \(R\) and all other coefficients
are nilpotent in \(R\).  Then apply the induction hypothesis to \(R\).  The
other parts are analogous.

\section*{Exercise 4}

\textbf{Issue.}  The proof is not merely incomplete; it proves a false
statement.  In general, \(A[x]\) does \emph{not} have Krull dimension \(0\),
and not every prime ideal of \(A[x]\) is maximal.  For example, if \(A=k\) is
a field, then
\[
(0)\subsetneq (x)
\]
are prime ideals of \(k[x]\).  The phrase ``Krull dimension of \(\fp\)'' is
also not standard; the relevant condition for \(\fp\) to be maximal would be
\(\dim(A/\fp)=0\).  The citation to Hartshorne, Proposition~1.7, gives
\(\dim\mathbb A^n=n\), not dimension \(0\).

The correct statement is
\[
J(A[x])=\Nil(A[x]),
\]
where \(J\) denotes the Jacobson radical and \(\Nil\) the nilradical.

\textbf{Correction.}  Let
\[
f=a_0+a_1x+\cdots+a_nx^n\in J(A[x]).
\]
For every \(g\in A[x]\), \(1-fg\) is a unit.  Taking \(g=x\), we see that
\(1-xf\) is a unit.  By Exercise~2(i), all coefficients of \(xf\) are
nilpotent; these coefficients are \(a_0,a_1,\ldots,a_n\).  Hence \(f\) is
nilpotent.  Therefore \(J(A[x])\subseteq\Nil(A[x])\).

Conversely, \(\Nil(R)\subseteq J(R)\) for every commutative ring \(R\): if
\(y^n=0\), then for every \(r\in R\),
\[
(1-yr)(1+yr+\cdots+(yr)^{n-1})=1-(yr)^n=1,
\]
so \(1-yr\) is a unit, which means \(y\in J(R)\).  Hence
\(\Nil(A[x])\subseteq J(A[x])\), and the two radicals are equal.

The remark that the same holds in
\(A[x_1,\ldots,x_n]\) is correct, by induction on \(n\).

\section*{Exercise 5}

\subsection*{(i)(b)}
The displayed inverse is missing the factor \(x^2\) in the third term.  The
beginning of the inverse series should read
\[
(a_0+a_1x+a_2x^2+\cdots)
\left(a_0^{-1}-a_1a_0^{-2}x+
(a_1^2a_0^{-3}-a_2a_0^{-2})x^2+\cdots\right)=1.
\]
More precisely, the coefficients \(b_n\) of the inverse are determined
recursively by
\[
b_0=a_0^{-1},\qquad
b_n=-a_0^{-1}\sum_{k=1}^{n}a_kb_{n-k}\quad(n\ge1).
\]

\subsection*{(ii)(b): the converse is false}
``No'' is correct, but a counterexample should be supplied.  Let
\[
A=\prod_{n\ge1}k[\varepsilon_n]/(\varepsilon_n^{n+1}),
\]
and let \(a_n\in A\) be the element whose \(n\)-th component is
\(\varepsilon_n\) and whose other components are \(0\).  Then
\(a_n^{n+1}=0\), so every \(a_n\) is nilpotent.  Consider
\[
f=\sum_{n\ge1}a_nx^n\in A[[x]].
\]
If \(f^N=0\), then in the \(n\)-th component we get
\(\varepsilon_n^N x^{nN}=0\).  Taking \(n>N\), this is impossible because
\(\varepsilon_n^N\neq0\) in \(k[\varepsilon_n]/(\varepsilon_n^{n+1})\).
Hence \(f\) is not nilpotent, although all of its coefficients are nilpotent.

\subsection*{(iii)(b): the classification of maximal ideals}

\textbf{Issue.}  The assertion
\[
\text{``the constant terms of an arbitrary ideal }\mathfrak I
\text{ form an ideal }I\text{ and }\mathfrak I_I=\mathfrak I\text{''}
\]
is false.  For example, in \(A[[x]]\), the ideal \(\mathfrak I=(x^2)\) has
constant-term ideal \(0\), but
\[
\mathfrak I_0=\{b_0+b_1x+\cdots\mid b_0=0\}=(x),
\]
which is strictly larger than \((x^2)\).

\textbf{Correction.}  Let \(\fm\) be a maximal ideal of \(A[[x]]\), and let
\[
\pi:A[[x]]\longrightarrow A,\qquad \pi\Bigl(\sum a_nx^n\Bigr)=a_0
\]
be the constant-term map, with kernel \((x)\).  First we show that
\(x\in\fm\).  If \(x\notin\fm\), then \(x\) is a unit modulo \(\fm\), so there
exists \(g\in A[[x]]\) with \(1-xg\in\fm\).  But \(1-xg\) has constant term
\(1\), hence is a unit in \(A[[x]]\), so \(\fm=A[[x]]\), contradiction.
Thus \(x\in\fm\).

Since \(x\in\fm\), we have
\[
\fm=\pi^{-1}(\pi(\fm)).
\]
Indeed, if \(\pi(f)\in\pi(\fm)\), then \(f-\pi(f)\in(x)\subseteq\fm\), so
\(f\in\fm\).  Therefore
\[
A[[x]]/\fm\cong A/\pi(\fm).
\]
Since \(\fm\) is maximal, \(\pi(\fm)\) is a maximal ideal of \(A\).  Thus every
maximal ideal of \(A[[x]]\) is of the form
\[
\fm=\pi^{-1}(I)=I+(x)
=\{b_0+b_1x+\cdots\mid b_0\in I\}
\]
with \(I\) maximal in \(A\).  Conversely, if \(I\) is maximal in \(A\), then
\(I+(x)\) is maximal because its quotient is \(A/I\).  This also proves
part (iv): the contraction of \(\fm\) is \(\pi(\fm)=\fm\cap A\), and \(\fm\)
is generated by \(\fm^c\) and \(x\).

\subsection*{(iii)(a)}
With the classification above, the equivalence is immediate.  If
\(a_0\in J(A)\), then for every maximal ideal \(I\) of \(A\),
\(a_0\in I\), hence \(f\in I+(x)\), which is an arbitrary maximal ideal of
\(A[[x]]\).  Thus \(f\in J(A[[x]])\).  Conversely, if
\(f\in J(A[[x]])\), then \(f\in I+(x)\) for every maximal \(I\subseteq A\),
so \(a_0\in I\) for every maximal \(I\), i.e.\ \(a_0\in J(A)\).

\subsection*{(iv)}
This is proved in the correction to (iii)(b).

\subsection*{(v)}
The idea in the original text is correct but unfinished.  For a prime ideal
\(\fp\subseteq A\), the ideal
\[
\fp A[[x]]=\fp[[x]]=\{a_0+a_1x+\cdots\mid a_i\in\fp\}
\]
is prime, because
\[
A[[x]]/\fp[[x]]\cong(A/\fp)[[x]],
\]
and \((A/\fp)[[x]]\) is a domain.  Its contraction to \(A\) is
\(\fp[[x]]\cap A=\fp\).  Hence every prime ideal of \(A\) is the contraction
of a prime ideal of \(A[[x]]\).  (Alternatively, \(\fp+(x)\) is prime with the
same contraction.)

\section*{Exercise 6}

\textbf{Issue.}  The proof is unfinished and its logic is incorrect.  The
statement ``if the nilradical equals the Jacobson radical, then every prime
ideal contains some maximal ideal'' is false: a prime ideal can contain the
intersection of all maximal ideals without containing any of them, e.g.\
\((0)\) in \(k[x]\).  The claimed equivalence with ``every prime ideal is
maximal'' is also false in general.

\textbf{Correction.}  Let \(J\) be the Jacobson radical and \(N\) the
nilradical.  Always \(N\subseteq J\).  Suppose \(J\not\subseteq N\), and choose
\(x\in J\setminus N\).  Then the principal ideal \((x)\) is not contained in
\(N\), so by hypothesis \((x)\) contains a nonzero idempotent \(e\).  Since
\(x\in J\), we have \(e\in J\).  But \(e(1-e)=0\), and \(1-e\) is a unit
because \(e\in J\); hence \(e=0\), contradiction.  Therefore
\(J\subseteq N\), and so \(J=N\).

\section*{Exercise 7}

Correct.  Minor wording: ``\(A/\fp\) is a integral domain'' should be ``an
integral domain''.  The hypothesis \(n>1\) is used only to ensure that
\(n-1>0\); the argument works for every \(n\ge2\).

\section*{Exercise 8}

\textbf{Issue.}  The nilradical is the intersection of all prime ideals, but
it need not itself be a prime ideal.  Hence it is not an element of the set
of prime ideals and cannot be used as a minimal element without further
argument.  For example, in \(A=k[x,y]/(xy)\),
\[
(x)\quad\text{and}\quad(y)
\]
are minimal prime ideals, while
\[
\Nil(A)=(x)\cap(y)=(xy)=(0)
\]
is not prime.

\textbf{Correction.}  Use Zorn's lemma.  The set of prime ideals is nonempty
because \(A\neq0\) has a maximal ideal, and maximal ideals are prime.  Let
\(\{\fp_i\}_{i\in I}\) be a chain of prime ideals, and put
\(\fp=\bigcap_i\fp_i\).  We show that \(\fp\) is prime.  Let \(xy\in\fp\) and
\(x\notin\fp\).  Choose \(i_0\) with \(x\notin\fp_{i_0}\).  For any \(j\):
if \(\fp_{i_0}\subseteq\fp_j\), then \(xy\in\fp_j\) and
\(x\notin\fp_j\), so \(y\in\fp_j\); if \(\fp_j\subseteq\fp_{i_0}\), then
\(y\notin\fp_j\) would imply \(x\in\fp_j\subseteq\fp_{i_0}\), contradiction,
so again \(y\in\fp_j\).  Hence \(y\in\fp\), and \(\fp\) is prime.  Zorn's
lemma now gives a minimal prime ideal.

\section*{Exercise 9}

The sketch is correct, but two details should be mentioned.  If
\(\fa=\Rad(\fa)\), then Proposition~1.14 gives
\[
\fa=\bigcap_{\fp\supseteq\fa}\fp,
\]
an intersection of prime ideals.  Conversely, an intersection of prime ideals
is radical, since each prime is radical and an intersection of radical ideals
is radical.  For \(\fa=(1)\), the family of primes containing \(\fa\) is
empty, and the intersection over the empty family is \(A\) by convention.

\section*{Exercise 10}

The proof is essentially correct, but several justifications are missing.

\begin{enumerate}[label=(\roman*)]
  \item \((i)\Rightarrow(ii)\): If \(A\) has exactly one prime ideal
        \(\fp\), then \(\fp\) is the unique maximal ideal and
        \(\Nil(A)=\fp\).  Every non-unit lies in a maximal ideal, hence in
        \(\fp=\Nil(A)\), so it is nilpotent.  Every element outside \(\fp\)
        is a unit.  The original text invokes Corollary~1.5 but omits the
        crucial equality \(\Nil(A)=\fp\).

  \item \((ii)\Rightarrow(iii)\): If \(\bar a\neq0\) in \(A/\Nil(A)\), then
        \(a\notin\Nil(A)\), so \(a\) is not nilpotent.  By (ii), \(a\) is a
        unit, hence \(\bar a\) is a unit.  The phrase
        ``\(\exists a^{-1}\in A\setminus\mathfrak R\)'' should be replaced by
        ``\(a\) is a unit''.

  \item \((iii)\Rightarrow(i)\): If \(A/\Nil(A)\) is a field, then
        \(\Nil(A)\) is a maximal ideal.  Every prime ideal contains
        \(\Nil(A)\), and since \(\Nil(A)\) is maximal, the only prime ideal is
        \(\Nil(A)\) itself.  The original conclusion is correct, but one
        should identify the unique prime as \(\Nil(A)\).
\end{enumerate}

\section*{Exercise 11}

\subsection*{(i)}
The chain should start with the Boolean identity for \(2x\):
\[
(2x)^2=2x,
\]
while also
\[
(2x)^2=4x^2=4x.
\]
Hence \(2x=4x\), so \(2x=0\).  As written, ``\(2x=4x^2\)'' is not justified.

\subsection*{(ii)}
Correct.

\subsection*{(iii)}
The two-generator case is correct.  For the induction, let
\(I=(x_1,\ldots,x_n)\) and set
\[
z=1-\prod_{i=1}^n(1+x_i).
\]
In a Boolean ring, \(1+x_i\) is idempotent because
\((1+x_i)^2=1+2x_i+x_i^2=1+x_i\).  Hence
\(e=\prod_{i=1}^n(1+x_i)\) is idempotent, and \(z=1-e\) is idempotent.
Moreover, \(ex_i=0\), since the factor \(1+x_i\) times \(x_i\) is
\(x_i+x_i^2=2x_i=0\).  Thus
\[
zx_i=(1-e)x_i=x_i,
\]
so \(x_i\in(z)\) for every \(i\).  Therefore \((z)=I\).

\section*{Exercise 12}

Correct.  A zero-divisor cannot be a unit, and in a local ring every non-unit
lies in the unique maximal ideal.  If \(e\neq0,1\) is idempotent, then
\(e(1-e)=0\), so both \(e\) and \(1-e\) are zero-divisors, hence non-units,
hence both lie in the maximal ideal; then \(1=e+(1-e)\) lies in the maximal
ideal, contradiction.  Only the grammar ``both \(e\) and \(1-e\) is'' should
be corrected to ``are''.

\section*{Exercise 13}

\textbf{Issue.}  The proof is not valid as written.  The notation
\(F\in A[f(x_f)]\) is unclear: \(F\) should be a polynomial in new variables
\(y_f\), evaluated at \(y_f=f(x_{f})\).  The ``highest degree'' cancellation
must be done with the coefficient of the highest power of \(f_0\); since
\(f_0\) is monic, \(f_0(x_{f_0})^{\alpha_0}\) is not a zero-divisor in the
polynomial ring over the other variables.  The equation
\[
f(x_f)^{\alpha_0}(\text{rest})=0
\]
is not the correct formulation of ``the highest-degree part vanishes''.
Finally, the termination argument needs the fact that the original relation
has no constant term in the \(y_f\) variables (or a minimality assumption);
it is not spelled out.  The parenthetical remark about the characteristic of
\(K\) is unnecessary and misleading: the argument works in every
characteristic.

\textbf{Correction.}  Suppose for contradiction that \(\fa=(1)\).  Then there
are finitely many \(f_1,\ldots,f_r\in\Sigma\) and
\(g_1,\ldots,g_r\in A\) such that
\[
\sum_{i=1}^r g_i f_i(x_{f_i})=1.
\]
Let \(L\) be a field extension of \(K\) containing a root \(\alpha_i\) of each
\(f_i\); for instance, take a splitting field of the product
\(f_1\cdots f_r\) over \(K\).  Evaluate the identity at
\[
x_{f_i}\longmapsto\alpha_i\quad(1\le i\le r),\qquad
x_f\longmapsto0\quad(f\notin\{f_1,\ldots,f_r\}).
\]
Then every \(f_i(x_{f_i})\) becomes \(f_i(\alpha_i)=0\), so the left-hand
side is \(0\), while the right-hand side is \(1\), contradiction.  Hence
\(\fa\neq(1)\).

\section*{Exercise 14}

\subsection*{(i)}
Zorn's lemma applies, but one should verify that the union of a chain of
ideals in \(\Sigma\) again lies in \(\Sigma\).  This is immediate: if every
element of every ideal in the chain is a zero-divisor, then every element of
the union is a zero-divisor.

\subsection*{(ii): invalid proof that a maximal element is prime}

\textbf{Issue.}  From \(ab\in I\) and \(u\neq0\) with \((ab)u=0\), one gets
\(a(bu)=0\).  This only shows that \(a\) is a zero-divisor if \(bu\neq0\);
if \(bu=0\), it shows that \(b\) is a zero-divisor instead.  More
importantly, even if \(a\) is a zero-divisor, maximality of \(I\) among ideals
consisting of zero-divisors does \emph{not} imply \(a\in I\).  For example,
in \(A=k[x,y]/(xy)\), the ideal \(I=(x)\) is maximal in \(\Sigma\), but \(y\)
is a zero-divisor and \(y\notin I\).

\textbf{Correction.}  Let \(I\) be maximal in \(\Sigma\), and suppose
\(ab\in I\).  Assume \(a\notin I\) and \(b\notin I\).  By maximality,
\(I+(a)\notin\Sigma\) and \(I+(b)\notin\Sigma\).  Hence there exist
\(i+ra\in I+(a)\) and \(j+sb\in I+(b)\) which are not zero-divisors.  Their
product
\[
(i+ra)(j+sb)=ij+isb+jra+rsab
\]
lies in \(I\), because \(i,j,ab\in I\).  Thus the product is a zero-divisor.
But the product of two non-zero-divisors is a non-zero-divisor: if
\((i+ra)(j+sb)z=0\), then \((j+sb)z=0\), and then \(z=0\).  Contradiction.
Therefore \(a\in I\) or \(b\in I\), so \(I\) is prime.

\subsection*{(iii)}
Correct.  To justify that every zero-divisor lies in some maximal element of
\(\Sigma\): if \(x\) is a zero-divisor with \(xu=0\) and \(u\neq0\), then every
element \(rx\) of the principal ideal \((x)\) satisfies \((rx)u=0\), so
\((x)\in\Sigma\).  By Zorn's lemma, \((x)\) is contained in a maximal element
of \(\Sigma\).  Hence the union of all maximal elements of \(\Sigma\) is
exactly the set of zero-divisors.

\section*{Exercise 15}

The answer ``Trivial'' with references is not a proof.  Here are the standard
verifications.

\begin{enumerate}[label=(\roman*)]
  \item \(V(E)=V((E))\) by definition.  Also
        \(V(\fa)=V(\Rad(\fa))\): a prime ideal contains \(\fa\) if and only if
        it contains \(\Rad(\fa)\).

  \item \(V(0)=X\), since every prime ideal contains \(0\); and
        \(V(1)=\varnothing\), since no prime ideal contains \(1\).

  \item For a family \((E_i)_{i\in I}\),
        \[
        \fp\in V\Bigl(\bigcup_iE_i\Bigr)
        \iff E_i\subseteq\fp\ \text{for all }i
        \iff \fp\in\bigcap_iV(E_i).
        \]
        Hence \(V(\bigcup_iE_i)=\bigcap_iV(E_i)\).

  \item For ideals \(\fa,\fb\),
        \[
        \fa\fb\subseteq\fa\cap\fb\subseteq\fa,\fb,
        \]
        so
        \[
        V(\fa)\cup V(\fb)\subseteq V(\fa\cap\fb)\subseteq V(\fa\fb).
        \]
        Conversely, if \(\fp\supseteq\fa\fb\) is prime, then
        \(\fa\subseteq\fp\) or \(\fb\subseteq\fp\), so
        \(\fp\in V(\fa)\cup V(\fb)\).  Hence all three sets are equal.
\end{enumerate}

\section*{Exercise 16}

\begin{enumerate}[label=(\roman*)]
  \item
        \[
        \Spec\mathbb Z=\{(0)\}\cup\{(p)\mid p\text{ is prime}\}.
        \]
        The zero ideal is missing in the original.

  \item \(\Spec\mathbb R=\{(0)\}\).  Correct.

  \item
        \[
        \Spec\mathbb C[x]=\{(0)\}\cup\{(x-a)\mid a\in\mathbb C\}.
        \]
        The zero ideal is missing in the original.

  \item
        \[
        \Spec\mathbb R[x]=\{(0)\}\cup\{(p)\mid p\in\mathbb R[x]
        \text{ is monic irreducible of degree }\ge1\}.
        \]
        Over \(\mathbb R\), the nonconstant irreducibles are the linear
        polynomials and the quadratics with negative discriminant.  The zero
        ideal is missing in the original.

  \item A complete description is
        \[
        \begin{aligned}
        \Spec\mathbb Z[x]
        &=\{(0)\}
        \cup\{(p)\mid p\text{ prime}\}\\
        &\quad\cup\{(p,f)\mid p\text{ prime},\ f\in\mathbb Z[x]\text{ monic and
        irreducible mod }p\}\\
        &\quad\cup\{(g)\mid g\in\mathbb Z[x]\text{ primitive and irreducible
        over }\mathbb Q\}.
        \end{aligned}
        \]
        Here \((p)\) is prime but not maximal; the maximal ideals are exactly
        the ideals \((p,f)\).  The ideal \((0)\) is the generic point, the
        ideals \((p)\) and \((g)\) are height-one primes, and the
        ideals \((p,f)\) are the height-two (closed) points.
\end{enumerate}

\section*{Exercise 17}

\subsection*{Basis}
The sets \(X_f=X\setminus V(f)\) are open by definition.  If
\(U=X\setminus V(E)\) is an arbitrary open set, then
\[
U=X\setminus V(E)=X\setminus\bigcap_{f\in E}V(f)
=\bigcup_{f\in E}(X\setminus V(f))
=\bigcup_{f\in E}X_f.
\]
Hence the \(X_f\) form a basis.

\subsection*{(i)--(iv)}
\begin{enumerate}[label=(\roman*)]
  \item
        \[
        X_f\cap X_g
        =X\setminus(V(f)\cup V(g))
        =X\setminus V(fg)
        =X_{fg}.
        \]
  \item \(X_f=\varnothing\iff V(f)=X\iff f\in\Nil(A)\).
  \item \(X_f=X\iff V(f)=\varnothing\iff f\) is a unit.
  \item
        \[
        X_f=X_g\iff V(f)=V(g)
        \iff\Rad((f))=\Rad((g)).
        \]
\end{enumerate}

\subsection*{(v)}
The original argument is essentially correct, but its quantifier notation
is unclear.  A clean version is: an open cover of \(X\) may be refined to a
cover by basic open sets \(X_{f_i}\).  If \(\bigcup_iX_{f_i}=X\), then
\(\bigcap_iV(f_i)=\varnothing\), so the ideal \(\fa=(f_i)\) is not contained
in any prime ideal, hence \(\fa=(1)\).  Therefore there are finitely many
indices \(i_1,\ldots,i_n\) and elements \(g_{i_j}\in A\) with
\[
1=\sum_{j=1}^ng_{i_j}f_{i_j}.
\]
Then \(\bigcup_{j=1}^nX_{f_{i_j}}=X\).  Thus \(X\) is quasi-compact.

\subsection*{(vi): radical, not the ideal itself}

\textbf{Issue.}  The condition for \(X_f\) to be covered by the \(X_{f_i}\) is
\[
f\in\Rad((f_i)_{i\in I}),
\]
not \(f\in(f_i)_{i\in I}\).  The latter is strictly stronger and is false in
general.  For example, in \(A=k[x]\), take \(f=x\) and \(f_1=x^2\).  Then
\(X_f=X_{x^2}\), so \(X_f\) is covered by the single set \(X_{x^2}\), but
\(x\notin(x^2)\).

\textbf{Correction.}  Suppose \(X_f\subseteq\bigcup_{i\in I}X_{f_i}\).  Then
\[
V(f)\supseteq\bigcap_{i\in I}V(f_i)=V(\fa),
\]
where \(\fa=(f_i)\).  Hence \(f\in\Rad(\fa)\), so \(f^n\in\fa\) for some
\(n\).  Since \(\fa\) is generated by the \(f_i\), there is a finite subset
\(J\subseteq I\) with \(f^n\in(f_i)_{i\in J}\).  Then
\[
V(f)=V(f^n)\supseteq V\bigl((f_i)_{i\in J}\bigr)
=\bigcap_{i\in J}V(f_i),
\]
so the finitely many \(X_{f_i}\) with \(i\in J\) cover \(X_f\).  Thus \(X_f\)
is quasi-compact.

\subsection*{(vii): the compactness argument is wrong}

\textbf{Issue.}  The negation of ``\(X_0\) is a finite union of basic open
sets'' is not ``the index family \(\mathcal F\) is infinite''.  An infinite
family can contain a finite subcover; for example, \(X\) itself is covered by
all the \(X_f\), and it is quasi-compact.  The attempted contradiction
therefore does not work.  In addition, the quantifiers in the original only
range over countable sequences \(X_{f_1},X_{f_2},\ldots\), whereas an open
cover can be an arbitrary family.

\textbf{Correction.}  Let \(X_0\subseteq X\) be open and quasi-compact.  Since
the \(X_f\) form a basis, write
\[
X_0=\bigcup_{f\in\mathcal F}X_f
\]
for some family \(\mathcal F\).  This is an open cover of \(X_0\).  By
quasi-compactness, finitely many of these sets cover \(X_0\), so \(X_0\) is a
finite union of basic open sets.  Conversely, a finite union of quasi-compact
spaces is quasi-compact, and each \(X_f\) is quasi-compact by (vi).

\section*{Exercise 18}

\subsection*{(i)--(iii)}
Correct.  In (i), \(\{x\}\) is closed if and only if
\(V(\fp_x)=\{\fp_x\}\), which is equivalent to \(\fp_x\) being maximal.
In (iii),
\[
y\in\overline{\{x\}}
\iff \fp_y\in V(\fp_x)
\iff \fp_x\subseteq\fp_y.
\]

\subsection*{(iv): \(V(\fp_x)\) is not a neighbourhood of \(x\)}

\textbf{Issue.}  The set \(V(\fp_x)=\overline{\{x\}}\) is closed, but it is
not necessarily a neighbourhood of \(x\): it need not contain an open set
containing \(x\).  The notation ``\(V(x)\)'' is also undefined; the intended
set is \(V(\fp_x)\).  The sentence ``I'll interpret neighborhood as closed
neighborhood'' is not the standard definition and leads to an incorrect
proof.

\textbf{Correction.}  Let \(x\neq y\) and suppose both
\(y\in\overline{\{x\}}\) and \(x\in\overline{\{y\}}\).  By (iii),
\(\fp_x\subseteq\fp_y\) and \(\fp_y\subseteq\fp_x\), hence
\(\fp_x=\fp_y\), so \(x=y\), contradiction.  Therefore at least one of
\[
y\notin\overline{\{x\}},\qquad x\notin\overline{\{y\}}
\]
holds.  If \(y\notin\overline{\{x\}}\), then
\(X\setminus\overline{\{x\}}\) is an open neighbourhood of \(y\) which does
not contain \(x\).  Otherwise, \(X\setminus\overline{\{y\}}\) is an open
neighbourhood of \(x\) which does not contain \(y\).  This is precisely the
\(T_0\) property.

\section*{Exercise 19}

\textbf{Issue.}  The proof is not a proof of irreducibility.  In the category
of ideals under inclusion, the initial object is the zero ideal \((0)\), not
the nilradical.  The nilradical is the intersection of all prime ideals; it
need not be prime, and it is not an element of \(\Spec(A)\) unless it is
prime.  The expression ``\(\Nil(A)\subseteq\Spec(A)\)'' is a type error; the
correct condition is \(\Nil(A)\in\Spec(A)\), i.e.\ \(\Nil(A)\) is prime.
The ``connected tree'' metaphor does not establish the topological
irreducibility of \(\Spec(A)\).

\textbf{Correction.}  A topological space \(X\) is irreducible if and only if
every two nonempty open subsets intersect.  Since the \(X_f\) form a basis,
it is enough to check basic open sets.  Recall that
\[
X_f\neq\varnothing\iff f\notin\Nil(A).
\]
Suppose \(\Nil(A)\) is prime.  If \(X_f\) and \(X_g\) are nonempty, then
\(f,g\notin\Nil(A)\), so \(fg\notin\Nil(A)\) because \(\Nil(A)\) is prime.
Hence
\[
X_f\cap X_g=X_{fg}\neq\varnothing.
\]
Thus every two nonempty basic open sets intersect, so \(\Spec(A)\) is
irreducible.

Conversely, suppose \(\Spec(A)\) is irreducible.  If \(fg\in\Nil(A)\), then
\[
X_f\cap X_g=X_{fg}=\varnothing.
\]
By irreducibility, one of \(X_f,X_g\) is empty, so \(f\in\Nil(A)\) or
\(g\in\Nil(A)\).  Hence \(\Nil(A)\) is prime.

\section*{Exercise 20}

This exercise is skipped in the original.  It is not mathematically wrong, but
it is incomplete.  A brief outline:

\begin{enumerate}[label=(\roman*)]
  \item If \(Y\subseteq X\) is irreducible, then \(\overline Y\) is
        irreducible: if \(U,V\) are nonempty open subsets of \(\overline Y\),
        then \(U\cap Y\) and \(V\cap Y\) are nonempty open subsets of \(Y\),
        so they meet, and hence \(U\cap V\neq\varnothing\).

  \item Apply Zorn's lemma to the set of irreducible subspaces of \(X\)
        ordered by inclusion.  The union of a chain of irreducible subspaces
        is irreducible, so maximal elements exist.

  \item A maximal irreducible subspace is closed, because its closure is
        irreducible and contains it.  The maximal irreducible subspaces cover
        \(X\), since every singleton is irreducible and is contained in a
        maximal one.  In a Hausdorff space, every irreducible subspace is a
        singleton: two distinct points have disjoint open neighbourhoods.
        Hence the irreducible components of a Hausdorff space are its
        points.

  \item For \(X=\Spec(A)\), the irreducible components are the closed sets
        \(V(\fp)\) with \(\fp\) a minimal prime ideal.  Indeed,
        \(V(\fp)\cong\Spec(A/\fp)\) is irreducible when \(\fp\) is prime, and
        every irreducible closed subset is of the form \(V(\fp)\) with
        \(\fp\) prime.  Minimality corresponds to minimality of \(\fp\).
\end{enumerate}

\section*{Exercise 21}

\subsection*{(i)}
Correct, apart from the typo \(Y_{\phi_{(f)}}\) in the last line, which should
be \(Y_{\phi(f)}\).

\subsection*{(ii)}
The conclusion is correct, but the displayed equality
\[
\phi(\phi^*(\fq))=\fq
\]
is false in general.  What is true is
\[
\phi(\phi^{-1}(\fq))\subseteq\fq.
\]
The equivalence
\[
\fa\subseteq\phi^{-1}(\fq)\iff\phi(\fa)\subseteq\fq
\]
is correct.  For example, if \(\phi:k[x]\to k[x,y]\) is the inclusion and
\(\fq=(x,y)\), then \(\phi^{-1}(\fq)=(x)\) and
\(\phi((x))=(x)\neq(x,y)\).

\subsection*{(iii)}
Only the easy inclusion is proved in the original; the reverse inclusion is
missing.  Here is the full proof.

We first show \(\phi^*(V(\fb))\subseteq V(\fb^c)\).  If \(\fq\supseteq\fb\),
then
\[
\phi^{-1}(\fb)\subseteq\phi^{-1}(\fq)=\phi^*(\fq),
\]
so \(\phi^*(\fq)\in V(\fb^c)\).  Since \(V(\fb^c)\) is closed,
\[
\overline{\phi^*(V(\fb))}\subseteq V(\fb^c).
\]

For the reverse inclusion, let \(\fp\in V(\fb^c)\).  Suppose
\(\fp\notin\overline{\phi^*(V(\fb))}\).  Then there is a basic open
neighbourhood \(X_f\) of \(\fp\) disjoint from \(\phi^*(V(\fb))\), with
\(f\notin\fp\).  For every \(\fq\in V(\fb)\), this means
\(\phi^*(\fq)\notin X_f\), i.e.\
\[
f\in\phi^{-1}(\fq),
\]
so \(\phi(f)\in\fq\).  Thus \(\phi(f)\) lies in every prime ideal containing
\(\fb\), hence
\[
\phi(f)\in\Rad(\fb).
\]
Therefore
\[
f\in\phi^{-1}(\Rad(\fb))=\Rad(\phi^{-1}(\fb))=\Rad(\fb^c).
\]
Since \(\fb^c\subseteq\fp\) and \(\fp\) is prime, \(\Rad(\fb^c)\subseteq\fp\),
so \(f\in\fp\), contradiction.  Hence
\(\fp\in\overline{\phi^*(V(\fb))}\), and the reverse inclusion holds.

\subsection*{(iv)}
The original proof is too vague.  Here is a precise argument.

If \(\phi\) is surjective, then \(\phi^*\) is injective: if
\(\phi^*(\fq_1)=\phi^*(\fq_2)\), then
\[
\fq_1=\phi(\phi^{-1}(\fq_1))
=\phi(\phi^*(\fq_1))
=\phi(\phi^*(\fq_2))
=\phi(\phi^{-1}(\fq_2))=\fq_2.
\]
Its image is \(V(\ker\phi)\): on one hand,
\(\ker\phi\subseteq\phi^{-1}(\fq)\) for every \(\fq\); on the other hand,
if \(\fp\supseteq\ker\phi\), then \(\fp=\phi^{-1}(\phi(\fp))\) and
\(\phi(\fp)\) is prime in \(B\cong A/\ker\phi\), so
\(\fp=\phi^*(\phi(\fp))\).

Finally, \(\phi^*\) is a closed map onto \(V(\ker\phi)\): for a closed set
\(V(\fb)\subseteq Y\), the argument in (iii) together with surjectivity gives
\[
\phi^*(V(\fb))=V(\fb^c).
\]
Hence \(\phi^*\) is a continuous bijection which is closed, so it is a
homeomorphism onto \(V(\ker\phi)\).  In particular, taking
\(B=A/\fN\) and \(\phi:A\to A/\fN\) the quotient map gives
\(\Spec(A)\cong\Spec(A/\fN)\).

\subsection*{(v)}
The chain of equivalences in the original is confused.  The correct proof is
short.  By (iii) with \(\fb=0\),
\[
\overline{\phi^*(Y)}
=\overline{\phi^*(V(0))}
=V(0^c)
=V(\ker\phi).
\]
Therefore \(\phi^*(Y)\) is dense in \(X\) if and only if
\[
V(\ker\phi)=X
\iff \ker\phi\subseteq\Nil(A).
\]
This is the ``more precisely'' statement.

The condition
\[
\Rad(\phi^{-1}(\Nil(B)))=\Nil(A)
\]
is in fact equivalent to \(\ker\phi\subseteq\Nil(A)\), but this equivalence
needs proof.  Since
\(\ker\phi\subseteq\phi^{-1}(\Nil(B))\) always, and
\(\Nil(A)\subseteq\phi^{-1}(\Nil(B))\) always, one has
\[
\Rad(\phi^{-1}(\Nil(B)))=\Nil(A)
\iff \phi^{-1}(\Nil(B))\subseteq\Nil(A)
\iff \ker\phi\subseteq\Nil(A).
\]
The original final step only proves the trivial inclusion
\(\ker\phi\subseteq\phi^{-1}(\Nil(B))\), so it does not establish the
desired equivalence.

\subsection*{(vi)}
Correct: for \(\fq\in\Spec(C)\),
\[
(\psi\circ\phi)^*(\fq)
=(\psi\circ\phi)^{-1}(\fq)
=\phi^{-1}(\psi^{-1}(\fq))
=\phi^*(\psi^*(\fq)).
\]

\subsection*{(vii)}
This part is left blank in the original.  Here is a solution outline.

Let \(A\) be an integral domain with exactly one nonzero prime ideal \(\fp\),
and let \(K=\Frac(A)\).  Put
\[
B=(A/\fp)\times K,\qquad \phi:A\longrightarrow B,\qquad
\phi(x)=(\bar x,x).
\]
Since \(\fp\) is the only nonzero prime, \(\Spec(A)=\{(0),\fp\}\), and
\(A/\fp\) is a field.  The spectrum of the product \(B\) is the disjoint union
\[
\Spec(B)=\{\fq_1,\fq_2\},
\]
where
\[
\fq_1=(0)\times K,\qquad \fq_2=(A/\fp)\times(0).
\]
Then
\[
\phi^*(\fq_1)=\phi^{-1}(\fq_1)=\fp,\qquad
\phi^*(\fq_2)=\phi^{-1}(\fq_2)=(0).
\]
Thus \(\phi^*\) is bijective.

However, \(\Spec(B)\) is discrete, while \(\Spec(A)\) is not: in
\(\Spec(A)\), the set \(\{(0)\}\) is open, but \(\{\fp\}\) is not open.
Indeed, \(X_f=\{(0)\}\) for \(f\in\fp\setminus\{0\}\), whereas \(f\notin\fp\)
gives \(X_f=X\), and \(f=0\) gives \(X_f=\varnothing\).  Hence no basic open
set equals \(\{\fp\}\).  Therefore \(\phi^*\) is a continuous bijection which
is not open, so it is not a homeomorphism.

\section*{Exercise 22}

\subsection*{(a)}
Correct.  The \(X_k\) are clopen, and the map
\(\Spec(A_k)\to X_k\) is a homeomorphism; this should be stated explicitly.

\subsection*{(b)(I) \((ii)\Rightarrow(i)\)}
The idea is correct, but one should add that \(X_1\) and \(X_2\) are nonempty,
disjoint, open and closed, and that \(X_1\cup X_2=X\).  Hence \(X\) is
disconnected.

\subsection*{(b)(II) \((iii)\Rightarrow(ii)\)}
There is a typo: the hypothesis is \(e^2=e\), not \(e^2=2\).  If
\(e\neq0,1\) is idempotent, then
\[
A\longrightarrow eA\times(1-e)A,\qquad a\longmapsto(ea,(1-e)a)
\]
is an isomorphism.  Both factors are nonzero because \(e\neq0\) and
\(1-e\neq0\).

\subsection*{(b)(III) \((i)\Rightarrow(iii)\)}
The construction is essentially correct, but several details should be
justified.  If \(V(I)\cap V(J)=\varnothing\), then \(I+J=(1)\), so there are
\(a\in I\), \(b\in J\) with \(a+b=1\).  Since \(V(IJ)=X\), we have
\(IJ\subseteq\Nil(A)\), so \((ab)^n=0\) for some \(n\).  Define
\[
s_1=\sum_{k=0}^{n-1}\binom{2n}{k}a^{2n-k}b^k,\qquad
s_2=\sum_{k=n+1}^{2n}\binom{2n}{k}a^{2n-k}b^k.
\]
Then the missing middle term is
\(\binom{2n}{n}(ab)^n=0\), so
\[
s_1+s_2=(a+b)^{2n}=1.
\]
Also \(s_1s_2=0\), because every term of the product is divisible by
\((ab)^{n+1}=0\).  Hence \(s_1(1-s_1)=0\), so \(s_1\) is idempotent, and
\(s_2=1-s_1\).  The idempotents are nontrivial: \(s_1\in I\) and \(I\) is
proper, so \(s_1\neq1\); \(s_2\in J\) and \(J\) is proper, so \(s_2\neq1\);
and \(s_1=0\) would give \(s_2=1\in J\), contradiction.  Thus \(s_1\neq0,1\).

\section*{Exercise 23}

\subsection*{(i)}
Correct.  For every \(f\),
\[
V(f)\cap V(1-f)=V((f)+(1-f))=V(1)=\varnothing,
\]
and
\[
V(f)\cup V(1-f)=V((f)(1-f))=V(0)=X,
\]
since \(f(1-f)=0\).  Hence \(V(f)=X_{1-f}\) is open, so \(X_f\) is closed.
The restriction to nonzero non-unital \(f\) is unnecessary.

\subsection*{(ii)}
Correct.  In a Boolean ring, every finitely generated ideal is principal by
Exercise~11(iii), so
\[
\bigcap_{k=1}^nV(f_k)
=V\Bigl(\sum_{k=1}^n(f_k)\Bigr)
=V((f))
\]
for some \(f\), hence
\[
X_{f_1}\cup\cdots\cup X_{f_n}=X_f.
\]

\subsection*{(iii): the clopen-set argument is wrong}

\textbf{Issue.}  The claim that one can choose \(f\in\fp\) which lies in no
other prime of \(Y\) is not Proposition~1.11(i), and it is false in general
when \(Y\) is infinite.  Moreover, even if such an \(f\) exists, the fact that
\(V(f)\) is a clopen set containing \(\fp\) does not imply \(V(f)=Y\); there
can be many clopen sets containing \(\fp\).  The sentence ``by
\(f\in\fp,\fp\in Y\) we have \(V(f)=Y\)'' is unjustified.

\textbf{Correction.}  Let \(Y\subseteq X\) be clopen.  Since \(X\) is
quasi-compact by Exercise~17(v), and \(Y\) is closed, \(Y\) is
quasi-compact.  Since \(Y\) is open, it is a union of basic open sets:
\[
Y=\bigcup_{f\in\mathcal F}X_f.
\]
By quasi-compactness, finitely many of these cover \(Y\), say
\(X_{f_1},\ldots,X_{f_n}\).  By (ii), there is an \(f\in A\) with
\[
X_{f_1}\cup\cdots\cup X_{f_n}=X_f.
\]
Hence \(Y=X_f\).  Conversely, every \(X_f\) is clopen by (i).  Therefore the
clopen subsets of \(X\) are exactly the sets \(X_f\).

\subsection*{(iv): the Hausdorff proof is wrong}

\textbf{Issue.}  The sets
\[
V_1=X\setminus\{\fq\},\qquad V_2=X\setminus\{\fp\}
\]
are \emph{open}, not closed, because \(\{\fp\}\) and \(\{\fq\}\) are closed
(Exercise~18(i)).  The reference to Exercise~17(iv) should also be
Exercise~17(v).

\textbf{Correction.}  First, \(X\) is quasi-compact by Exercise~17(v).  To
prove that \(X\) is Hausdorff, let \(\fp\neq\fq\) be distinct points.  Since
they are maximal ideals (Exercise~11(ii)), \(\fp+\fq=A\), so there are
\(a\in\fp\), \(b\in\fq\) with \(a+b=1\).  In a Boolean ring, \(a\) and \(b\)
are idempotents, and \(b=1-a\).  Then
\[
X_a=\{ \fr\in X\mid a\notin\fr\},\qquad
X_b=\{ \fr\in X\mid b\notin\fr\}
\]
are disjoint clopen sets, with \(\fq\in X_a\) (because \(a\notin\fq\)) and
\(\fp\in X_b\) (because \(b\notin\fp\)).  Hence \(X\) is Hausdorff.

\section*{Exercise 24}

\subsection*{(a)(I)(iv): the associativity computation is wrong}

\textbf{Issue.}  The displayed computation
\[
(a+b)+c=(ab'\vee a'b)c\vee(ab+a'b')c
\]
does not correctly expand \((a+b)+c\), and the final expression is not
symmetric in \(a,b,c\).

\textbf{Correction.}  Recall \(x+y=(xy')\vee(x'y)\).  Then
\[
\begin{aligned}
(a+b)+c
&=\bigl((ab'\vee a'b)c'\bigr)\vee\bigl((ab'\vee a'b)'c\bigr)\\
&=(ab'c'\vee a'bc')\vee\bigl((a'\vee b)(a\vee b')c\bigr)\\
&=ab'c'\vee a'bc'\vee a'b'c\vee abc.
\end{aligned}
\]
This expression is symmetric in \(a,b,c\) (it is the set of terms with an
odd number of unprimed variables), so addition is associative.

\subsection*{(a)(II)(iii)}
``Associativity of multiplication is given by the associativity of
\(\vee\)'' should be ``of \(\wedge\)'', since multiplication is \(a\wedge b\).

\subsection*{(b)(iii): the absorption computation is wrong}

\textbf{Issue.}  The line
\[
a\vee(a\wedge b)=a+ab+aab=3a
\]
is incorrect; \(a+ab+aab=a+2ab\), not \(3a\).

\textbf{Correction.}  Since \(a^2=a\) and \(2ab=0\),
\[
a\vee(a\wedge b)
=a\vee(ab)
=a+ab+a^2b
=a+ab+ab
=a+2ab
=a.
\]
The other absorption law follows by duality.

\subsection*{Missing details in (b)}
To show that \(A\) is a Boolean lattice with respect to \(a\le b\iff a=ab\),
one should also verify that \(\le\) is a partial order (reflexivity,
antisymmetry, transitivity), that \(a\wedge a'=0\) (in addition to
\(a\vee a'=1\)), and that complements are unique.  These are short
verifications.  For example,
\[
a\wedge a'=a(1-a)=a-a^2=0,
\]
and if \(b\) is another complement of \(a\), then
\[
b=b\wedge1=b\wedge(a\vee a')=(b\wedge a)\vee(b\wedge a')
=0\vee(b\wedge a')=b\wedge a',
\]
so \(b\le a'\); similarly \(a'\le b\), hence \(b=a'\).

\section*{Exercise 25}

The statement is the conclusion of Stone's theorem, not a proof.  The
intended argument is:

\begin{enumerate}[label=(\roman*)]
  \item By Exercise~24, \(A(L)\) is a Boolean ring.
  \item By Exercise~23, the clopen subsets of \(X=\Spec(A(L))\) are exactly
        the sets \(X_f\) with \(f\in A(L)\).
  \item The map
        \[
        L\longrightarrow\{\text{clopen subsets of }X\},\qquad
        a\longmapsto X_a
        \]
        preserves the Boolean lattice structure (order, sup, inf,
        complement).
  \item This map is bijective by Exercise~23(iii).
\end{enumerate}
Hence \(L\) is isomorphic to the lattice of clopen subsets of the compact
Hausdorff space \(X\).

\section*{Exercise 26}

\textbf{Issue.}  In the first inclusion, the implication
\[
\fm_y\in\widetilde U_f,\ f\in\fm_y\ \Longrightarrow\ f(y)\neq0
\]
has the wrong sign.  By definition,
\[
\widetilde U_f=\{\fm\in\Max(C(X))\mid f\notin\fm\}.
\]
Thus \(\fm_y\in\widetilde U_f\) means \(f\notin\fm_y\), i.e.\ \(f(y)\neq0\).
Also, the second inclusion is written as \(U_f\subseteq\widetilde U_f\), but
these are subsets of different spaces; it should be
\(\mu(U_f)\subseteq\widetilde U_f\), or equivalently
\(U_f\subseteq\mu^{-1}(\widetilde U_f)\).

\textbf{Correction.}  Let \(\fm\in\widetilde U_f\).  By (i), \(\fm=\fm_y\) for
some \(y\in X\).  Since \(\fm_y\in\widetilde U_f\), we have
\(f\notin\fm_y\), hence \(f(y)\neq0\), so \(y\in U_f\) and
\(\fm=\mu(y)\in\mu(U_f)\).  Thus
\(\widetilde U_f\subseteq\mu(U_f)\).

Conversely, if \(x\in U_f\), then \(f(x)\neq0\), so \(f\notin\fm_x\), hence
\(\mu(x)=\fm_x\in\widetilde U_f\).  Thus
\(\mu(U_f)\subseteq\widetilde U_f\).

Therefore \(\mu(U_f)=\widetilde U_f\).  Since \(\mu\) is bijective by (i) and
(ii), and the sets \(U_f\) (respectively \(\widetilde U_f\)) form a basis of
the topology, \(\mu\) is a homeomorphism.

\section*{Exercise 27}

\textbf{Issue.}  The argument is circular.  The step
\[
\text{``}\exists x\in X\text{ such that }\varphi^{-1}(\fm)
\text{ is the maximal vanishing ideal of }x\text{''}
\]
is precisely the weak Nullstellensatz (every maximal ideal of
\(k[t_1,\ldots,t_n]\) is of the form
\((t_1-x_1,\ldots,t_n-x_n)\)).  This is not available at this point in the
book; the text itself says that surjectivity of \(\mu\) is one form of
Hilbert's Nullstellensatz, proved in Chapter~7.  There is also a notation
error: \(\varphi:X\to X/I(X)\) should be
\(\varphi:k[t_1,\ldots,t_n]\to P(X)\).

\textbf{Correction.}  The correct statement at this point is the equivalence:
\(\mu\) is surjective if and only if every maximal ideal of \(P(X)\) is a
vanishing ideal \(\fm_x\).  Under the quotient map
\(\varphi:k[t_1,\ldots,t_n]\to P(X)\), maximal ideals of \(P(X)\) correspond
to maximal ideals of \(k[t_1,\ldots,t_n]\) containing \(I(X)\).  Thus
surjectivity of \(\mu\) is equivalent to the weak Nullstellensatz.  A proof
should either cite Chapter~7 or prove the Nullstellensatz; the argument in
the original file assumes it.

\section*{Exercise 28}

\textbf{Issue.}  The proof omits the crucial verification that the map
\[
\varphi=(F_1,\ldots,F_m):X\longrightarrow k^m
\]
actually lands in \(Y\).  It is not enough to say that
\((F_1,\ldots,F_m)\) ``defines a regular mapping \(X\to Y\)''; one must show
that for every \(x\in X\), the point \((F_1(x),\ldots,F_m(x))\) satisfies all
the equations defining \(Y\).  The proof also does not prove injectivity of
the correspondence.

\textbf{Correction.}  Let \(\Phi:P(Y)\to P(X)\) be a \(k\)-algebra
homomorphism, and let \(\xi_i\) be the coordinate functions on \(Y\).  Put
\[
F_i=\Phi(\xi_i)\in P(X)\qquad(1\le i\le m).
\]
For \(g\in I(Y)\), we have \(g(\xi_1,\ldots,\xi_m)=0\) in \(P(Y)\), hence
\[
0=\Phi(g(\xi_1,\ldots,\xi_m))=g(F_1,\ldots,F_m)
\]
in \(P(X)\).  Therefore, for every \(x\in X\),
\[
g(F_1(x),\ldots,F_m(x))=0.
\]
This shows that \((F_1(x),\ldots,F_m(x))\in Y\), so
\(\varphi=(F_1,\ldots,F_m)\) is a regular map \(X\to Y\).

Let \(\varphi^*:P(Y)\to P(X)\) be the induced homomorphism
\(\eta\mapsto\eta\circ\varphi\).  On the generators \(\xi_i\),
\[
\varphi^*(\xi_i)=\xi_i\circ\varphi=F_i=\Phi(\xi_i).
\]
Since \(P(Y)\) is generated as a \(k\)-algebra by
\(\xi_1,\ldots,\xi_m\), it follows that \(\varphi^*=\Phi\).  Hence the
correspondence is surjective.

For injectivity, suppose \(\varphi,\psi:X\to Y\) induce the same
\(k\)-algebra homomorphism.  Then for every \(i\),
\[
\xi_i\circ\varphi=\xi_i\circ\psi,
\]
so the \(i\)-th coordinates of \(\varphi(x)\) and \(\psi(x)\) agree for every
\(x\in X\).  Hence \(\varphi=\psi\).  Therefore the correspondence is
bijective.

\section*{Unfinished or sketchy exercises}

The following exercises are not necessarily wrong, but they are incomplete:
\begin{itemize}[leftmargin=1.5em]
  \item \textbf{Ex.~3}: the induction is only indicated.
  \item \textbf{Ex.~5(v)}: the idea is correct; see the correction above.
  \item \textbf{Ex.~9}: the proof is only a reference to Proposition~1.14.
  \item \textbf{Ex.~15}: ``Trivial'' with references is not a proof.
  \item \textbf{Ex.~16(v)}: \(\Spec\mathbb Z[x]\) is left unfinished.
  \item \textbf{Ex.~17}: the basis and parts (i)--(iv) are asserted without
        proof; (vi) and (vii) contain errors.
  \item \textbf{Ex.~20}: skipped.
  \item \textbf{Ex.~21(iii),(v),(vii)}: (iii) and (v) are incomplete;
        (vii) is blank.
  \item \textbf{Ex.~25}: only the statement of Stone's theorem is given.
  \item \textbf{Ex.~27}: the argument is conditional on the
        Nullstellensatz.
  \item \textbf{Ex.~28}: the key verification and injectivity are missing.
\end{itemize}

\section*{Minor typographical corrections}

\begin{itemize}[leftmargin=1.5em]
  \item Ex.~1: replace \(1-x^r\) by \(1-(-x)^r\) in the intermediate identity.
  \item Ex.~2(i)(b): \(a_ns^n\) should be \(a_nx^n\).
  \item Ex.~2, remark: ``\(n+m-1\) anti-diagonals'' should be
        ``\(n+m+1\) anti-diagonals''.
  \item Ex.~5(i)(b): the third term of the inverse series should contain
        \(x^2\).
  \item Ex.~10: the nilradical is usually denoted \(\mathfrak N\), not
        \(\mathfrak R\); ``a integral domain'' should be ``an integral
        domain''.
  \item Ex.~11(i): the equality \(2x=4x^2\) needs the justification
        \((2x)^2=2x\).
  \item Ex.~21(i): \(Y_{\phi_{(f)}}\) should be \(Y_{\phi(f)}\).
  \item Ex.~21(ii): \(\phi(\phi^{-1}(\fq))=\fq\) should be
        \(\phi(\phi^{-1}(\fq))\subseteq\fq\).
  \item Ex.~22(b)(II): \(e^2=2\) should be \(e^2=e\).
  \item Ex.~23(iv): ``Ex-17(iv)'' should be ``Ex-17(v)''.
  \item Ex.~24(a)(II)(iii): ``associativity of \(\vee\)'' should be
        ``associativity of \(\wedge\)''.
  \item Ex.~24(b)(iii): ``\(=3a\)'' should be ``\(=a+2ab=a\)''.
  \item Ex.~26: \(f\in\fm_y\) should be \(f\notin\fm_y\); and
        \(U_f\subseteq\widetilde U_f\) should be
        \(\mu(U_f)\subseteq\widetilde U_f\).
  \item Ex.~27: \(\varphi:X\to X/I(X)\) should be
        \(\varphi:k[t_1,\ldots,t_n]\to P(X)\).
  \item Ex.~28: the diagram should have
        \(P(Y)=k[t_1,\ldots,t_m]/I(Y)\); ``For any \(k\)-algebra
        homomorphism \(\phi\)'' should be ``\(\Phi\)''; and ``regualr''
        should be ``regular''.
\end{itemize}

\end{document}
